\section{Repository Artifact Note}

This paper is included as the final artifact in the repository's linked-paper upgrade example. The manuscript was generated from a user-provided source paper that the user identified as their own work: \citet{corcoran2018spatial}. The upgrade path taken here is a \emph{pivot}, not a cosmetic rewrite.

The full artifact chain for this example lives one directory above this paper and includes:
\begin{itemize}[leftmargin=1.4em]
  \item source-paper notes and review,
  \item a literature update,
  \item an explicit publication decision,
  \item a breakthrough plan,
  \item an improvement diff,
  \item a narrative report for paper generation, and
  \item this compiled manuscript plus an improvement log.
\end{itemize}

The package also contains a reproducibility bundle:
\begin{itemize}[leftmargin=1.4em]
  \item build and validation scripts in \texttt{../code/},
  \item a source-data manifest in \texttt{../data/},
  \item dedicated high-resolution figure assets in \texttt{../figure\_assets/}, and
  \item a scored review in \texttt{../review/}.
\end{itemize}

The manuscript is intentionally honest about scope. It combines a perspective-style reinterpretation with a deterministic synthetic serialization probe. For that reason, the executable code in this package covers probe generation, table and figure generation, paper building, and artifact validation rather than model training on a real biomolecular benchmark.
